% Advisor-facing method outline for ToolFrame 6.4.
\documentclass[journal]{IEEEtran}

\usepackage{amsmath,amssymb}
\usepackage{booktabs}
\usepackage{xcolor}

\newcommand{\todo}[1]{\textcolor{red}{TODO: #1}}
\newcommand{\method}{ToolFrame}

\title{\method{}: Let the Tool Define the Frame for Dexterous Manipulation}
\author{Advisor Brief Draft}

\begin{document}
\maketitle

\section*{Title Candidates}
\textbf{Preferred:}
\begin{quote}
\method{}: Let the Tool Define the Frame for Dexterous Manipulation
\end{quote}

\textbf{Safer RA-L-style alternative:}
\begin{quote}
\method{}: Tool-Centric Visual Grounding for Dexterous Manipulation
\end{quote}

\textbf{Method-oriented alternative:}
\begin{quote}
\method{}: Tool-Centric Compilation of Visual Evidence for Dexterous Tool Use
\end{quote}

\section{Core Thesis}
Traditional VLM tool-use pipelines often treat the VLM as a planner that emits
actions, trajectories, or waypoints. \method{} instead makes the tool, rather
than the hand, the organizing frame for dexterous manipulation:
\begin{quote}
The tool defines the frame; the hand follows.
\end{quote}

The method decomposes tool use into three layers:
\begin{enumerate}
  \item \textbf{Functional tool geometry}: functional part, contact point, and
  motion axis.
  \item \textbf{Scene action evidence}: RGB-D start point, path points, end
  point, and direction.
  \item \textbf{Whole-tool execution targets}: compiled whole-tool $SE(3)$
  subgoals executed by a dexterous hand policy.
\end{enumerate}

\section{Abstract Skeleton}
\begin{enumerate}
  \item \textbf{Background}: Dexterous tool use requires binding functional
  tool geometry to task-relevant scene evidence.
  \item \textbf{Problem}: Direct VLM trajectory generation entangles semantic
  grounding, metric geometry, control feasibility, and hidden benchmark
  assumptions.
  \item \textbf{Method}: \method{} restricts VLMs to sparse public evidence
  selection: tool refs and scene refs.
  \item \textbf{Compiler}: A geometry-constrained compiler binds
  \texttt{tool.contact\_point} and \texttt{tool.motion\_axis} to lifted RGB-D
  scene refs and outputs whole-tool $SE(3)$ subgoals.
  \item \textbf{Execution}: A SimToolReal fixed-goal dexterous policy executes
  these subgoals.
  \item \textbf{Evaluation}: Task contracts evaluate final success and
  attribute failures by module.
  \item \textbf{Results}: \todo{Fill final results from frozen
  launcher\_summary files and per-task summary.json files.}
\end{enumerate}

\section{Introduction}
\subsection{Motivation}
Tool use is functional geometric grounding. A brush provides broad lateral
contact, an eraser provides a flat wiping patch, a hammer provides an impact
face and strike axis, a screwdriver provides a tip and shaft axis, and a marker
provides a point tip and trace path. The dexterous hand should not merely move
an end-effector; it should bring the functional part of the tool to the scene
with the right pose and direction.

\subsection{Problem}
VLMs can identify task-relevant visual evidence, but asking them to directly
generate robot trajectories mixes semantic grounding, metric geometry, policy
feasibility, and benchmark hidden state. When execution fails, it is unclear
whether the error came from tool grounding, scene grounding, geometric
compilation, policy execution, or the task effect.

\subsection{Key Idea}
\method{} treats the VLM as an evidence selector:
\begin{itemize}
  \item VLMs select sparse scene refs.
  \item VLM / part evidence selects functional tool refs.
  \item The compiler produces metric whole-tool $SE(3)$ subgoals.
  \item A dexterous policy executes fixed goals.
  \item Task contracts evaluate success and failure attribution.
\end{itemize}

\subsection{Contributions}
\begin{enumerate}
  \item A tool-centric visual grounding interface using
  \texttt{tool.contact\_point} and \texttt{tool.motion\_axis}.
  \item A geometry-constrained compiler that binds tool refs and RGB-D scene
  refs into whole-tool $SE(3)$ subgoals.
  \item A closed-loop pipeline from public visual evidence to SimToolReal
  fixed-goal dexterous rollout.
  \item A task-contract evaluation protocol for final success and module-level
  failure attribution.
\end{enumerate}

\section{Related Work}
\subsection{VLMs for Robot Planning}
Discuss VLM / LLM robot task planning, visual grounding, code-as-policies, and
language-conditioned manipulation. Difference: \method{} does not use the VLM
as a trajectory generator.

\subsection{Tool Affordance and Functional Part Grounding}
Discuss affordance learning, part segmentation, functional region prediction,
and contact point estimation. Difference: \method{} uses functional tool refs
as compiler inputs, not only as semantic labels.

\subsection{Constraint-Based Manipulation and Subgoal Planning}
Discuss task-space constraints, contact-rich manipulation, keyframe planning,
and $SE(3)$ pose optimization. Difference: \method{} compiles sparse visual
refs into whole-tool $SE(3)$ subgoals.

\subsection{Tool-Use Benchmarks and Evaluation}
Discuss manipulation benchmarks, tool-use benchmarks, and success metrics.
Difference: \method{} separates task success from module-level failure
attribution.

\section{Problem Formulation and Evidence Boundary}
\subsection{Allowed Inputs}
The method may use task language, public RGB-D scene evidence, camera
intrinsics/extrinsics, public tool mesh/part evidence, reviewed-or-pending
public tool refs, and robot/gripper state needed for execution.

\subsection{Forbidden Inputs}
The method must not use hidden evaluator target refs, oracle target axes,
oracle trajectories, simulator target centers for snapping refs, tool-family
fallback, task-family fallback, primitive fallback, repair-loop fallback, or
seed-specific shortcuts.

\subsection{Reference Representation}
Tool side:
\begin{quote}
\texttt{tool.contact\_point}, \texttt{tool.motion\_axis}
\end{quote}

Scene side:
\begin{quote}
\texttt{scene.start\_point}, \texttt{scene.path\_points},
\texttt{scene.end\_point}, \texttt{scene.direction}
\end{quote}

Core objective: find whole-tool $SE(3)$ poses such that the selected tool
contact point reaches the selected scene refs while the selected tool motion
axis aligns with the scene direction.

\section{Method: ToolFrame}
\subsection{Overview}
Figure~1 should show four modules:
\begin{quote}
(A) Tool-Centric Prior Extraction;
(B) Scene Key-Points Grounding;
(C) Geometry-Constrained Optimizer;
(D) Policy Rollout and Evaluation.
\end{quote}

Suggested caption:
\begin{quote}
\method{} overview. The tool defines the manipulation frame: public tool-part
evidence yields a functional part, contact point, and motion axis; public
RGB-D scene evidence yields sparse scene references; a tool-centric $SE(3)$
compiler binds them into whole-tool subgoals executed by a fixed-goal
dexterous policy. The VLM selects evidence only and does not output robot
trajectories.
\end{quote}

\subsection{A. Tool-Centric Prior Extraction}
\textbf{Inputs}: tool mesh, segmented tool-part evidence, task language, and
reviewed visual evidence.

\textbf{Process}: generate candidate functional parts; select the functional
part from public evidence; select contact point and motion axis on the chosen
part; convert the selection to the tool frame.

\textbf{Outputs}: \texttt{tool.contact\_point}, \texttt{tool.motion\_axis},
\texttt{tool\_part\_graph}, part-selection evidence, and axis-binding
evidence.

\subsection{B. Scene Key-Points Grounding}
\textbf{Inputs}: public RGB-D image, depth, camera calibration, and natural
task language.

\textbf{Process}: VLM selects start point, optional path points, end point, and
scene direction; RGB-D lifting converts 2D refs into 3D refs.

\textbf{Outputs}: \texttt{scene.start\_point\_3d},
\texttt{scene.path\_points\_3d}, \texttt{scene.end\_point\_3d}, and
\texttt{scene.direction\_3d}.

\subsection{C. Geometry-Constrained Optimizer}
\textbf{Inputs}: \texttt{tool.contact\_point}, \texttt{tool.motion\_axis},
lifted scene refs, and task-stage constraints.

\textbf{Constraints}: point match, point offset, axis alignment, move along,
orientation offset about axis, and explicit diagnostic clearance.

\textbf{Outputs}: \texttt{tool\_pose\_subgoals.json},
\texttt{simtoolreal\_fixed\_goals.json}, and compiler report.

This section should emphasize that the compiler does not infer task primitives
such as sweep, strike, wipe, or insert. It only solves geometric constraints
over tool and scene refs.

\subsection{D. Policy Rollout and Evaluation}
\textbf{Inputs}: fixed whole-tool subgoals, SimToolReal scene state, and
rollout config.

\textbf{Process}: the fixed-goal dexterous policy follows subgoals; rollout
produces videos, contact evidence, and tracking diagnostics; task contracts
evaluate final success; failure attribution records the first failed module.

\textbf{Outputs}: \texttt{rollout\_summary.json},
\texttt{visual\_evidence/manifest.json}, task success checks, failure
attribution, and \texttt{summary.json}.

\section{Benchmark and Tasks}
\subsection{Tool Coverage}
\begin{table}[h]
\centering
\caption{Tool functional geometries targeted by \method{}.}
\begin{tabular}{ll}
\toprule
Tool & Functional geometry \\
\midrule
Brush & broad lateral contact and sweep axis \\
Eraser & flat bottom contact patch and wipe axis \\
Hammer & impact face and strike axis \\
Screwdriver & tip and shaft/insertion axis \\
Marker & point tip and writing/contact axis \\
\bottomrule
\end{tabular}
\end{table}

\subsection{Scene Coverage}
\begin{table}[h]
\centering
\caption{Scene families.}
\begin{tabular}{ll}
\toprule
Scene & Purpose \\
\midrule
Tabletop & horizontal tabletop manipulation \\
Side-board & vertical board manipulation \\
\bottomrule
\end{tabular}
\end{table}

\subsection{Metrics}
Main metrics:
\begin{quote}
functional grounding rate, scene grounding rate, subgoal validity rate, policy
execution rate, task success rate, and failure attribution.
\end{quote}
Task success rate is the final task metric; other metrics explain where the
pipeline fails.

\section{Experiments}
\subsection{Main Result Table}
Final tables should report:
\begin{quote}
Task, Tool, Scene, N, Functional, Scene, Subgoal, Policy, Success, Failure.
\end{quote}
Numbers should only come from frozen \texttt{launcher\_summary} files and
matching per-task \texttt{summary.json} files.

\subsection{Qualitative Rollout}
Show representative rollouts for brush sweep, eraser wipe, hammer strike,
screwdriver insert, and marker trace. Each example should show:
\begin{quote}
tool refs $\rightarrow$ scene refs $\rightarrow$ compiled subgoals
$\rightarrow$ rollout result.
\end{quote}

\subsection{Failure Attribution}
Show whether failures occur in functional grounding, scene grounding, subgoal
validity, policy execution, task effect, or needs-review cases.

\section{Figures and Tables}
\begin{itemize}
  \item Figure 1: overall pipeline.
  \item Figure 2: five tools and functional refs.
  \item Figure 3: tabletop and side-board scene grounding.
  \item Figure 4: whole-tool $SE(3)$ subgoals and dexterous rollout.
  \item Figure 5: failure attribution or success summary.
  \item Table 1: tool functional geometry coverage.
  \item Table 2: task, scene, and contract status.
  \item Table 3: main result table.
\end{itemize}

\section{Limitations}
\begin{itemize}
  \item Requires visible RGB-D evidence and accurate lifting.
  \item Tool refs require public part evidence and human review.
  \item Fixed-goal policy may fail even with valid compiled subgoals.
  \item Side-board and marker-trace tasks may require stronger contracts.
  \item No online repair loop or seed-specific shortcut is used.
\end{itemize}

\section{Conclusion}
\method{} structures dexterous tool use as tool-centric evidence grounding and
geometric subgoal compilation: the tool defines the frame, and the hand
follows.

\end{document}
